% GENERATED FILE. Edit canonical lessons, then run npm run generate:books.
% canonical-source-hash: fnv1a64:3300f7c9f4101738
% canonical-lessons: PT-C527-realizador, PT-C527-filmar, PT-C527-guiao, PT-C527-documentario, PT-C527-elenco

\chapter{Atrás das câmaras}
\label{ch:pt-atras-das-camaras}
\hlchaptermodality{527}

\begin{chapteropening}
\textbf{By the end of this chapter:} \emph{I can talk about directors, scripts and documentaries.}

Complete the last lesson of chapter 527: I can talk about directors, scripts and documentaries.
\end{chapteropening}

\section[o realizador]{o realizador — (film) director}

\label{lesson:PT-C527-realizador}



\begin{quote}
Before the new one: say the Portuguese for subtitled, then the Portuguese for dubbing.
\end{quote}

\subsection*{You'll want to know: o realizador}
\textbf{o realizador} — \textquotedblleft{}(film) director\textquotedblright{}.

\begin{quote}
O realizador ganhou um prémio pelo filme.
\end{quote}

\begin{quote}
\emph{The director won a prize for the film.}
\end{quote}

In Brazil: \emph{o diretor}.

\begin{cousinweb}
\emph{realizador} is built on \emph{realizar} \textquotedblleft{}to carry out, to make\textquotedblright{}, which is usually traced to \emph{real} \textquotedblleft{}real\textquotedblright{} and so to Latin \emph{rēs} \textquotedblleft{}a thing\textquotedblright{}.
\end{cousinweb}

\begin{grammarlens}[title={What you've built}]
One more word for telling, explaining and describing at B1.
\end{grammarlens}

\subsection*{Your turn}
\begin{itemize}
\raggedright
  \item \emph{Say it:} \emph{o realizador}
  \item \emph{Say it:} \emph{o realizador}, once more
  \item \emph{Say it:} a sentence of your own with \emph{o realizador}
  \item \emph{Recall:} say the Portuguese for subtitled, then the Portuguese for dubbing, then say \emph{o realizador} again
\end{itemize}

\begin{tcolorbox}[breakable,colback=teal!4,colframe=teal!35!black,title={Before you move on}]
What does o realizador mean? (\textquotedblleft{}(film) director\textquotedblright{}.) Say it once more.
\end{tcolorbox}
\section[filmar]{filmar — to film}

\label{lesson:PT-C527-filmar}



\begin{quote}
Before the new one: say the Portuguese for dubbing, then the Portuguese for (film) director.
\end{quote}

\subsection*{You'll want to know: filmar}
\textbf{filmar} — \textquotedblleft{}to film\textquotedblright{}.

\begin{quote}
Filmaram o casamento com o telemóvel.
\end{quote}

\begin{quote}
\emph{They filmed the wedding on their mobile.}
\end{quote}

Built on \emph{o filme}.

\begin{cousinweb}
\emph{filmar} is built on \emph{filme}, which was borrowed from English \emph{film} \textquotedblleft{}a thin skin\textquotedblright{}.
\end{cousinweb}

\begin{grammarlens}[title={What you've built}]
One more word for telling, explaining and describing at B1.
\end{grammarlens}

\subsection*{Your turn}
\begin{itemize}
\raggedright
  \item \emph{Say it:} \emph{filmar}
  \item \emph{Say it:} \emph{filmar}, once more
  \item \emph{Say it:} a sentence of your own with \emph{filmar}
  \item \emph{Recall:} say the Portuguese for dubbing, then the Portuguese for (film) director, then say \emph{filmar} again
\end{itemize}

\begin{tcolorbox}[breakable,colback=teal!4,colframe=teal!35!black,title={Before you move on}]
What does filmar mean? (\textquotedblleft{}To film\textquotedblright{}.) Say it once more.
\end{tcolorbox}
\section[o guião]{o guião — script, screenplay}

\label{lesson:PT-C527-guiao}



\begin{quote}
Before the new one: say the Portuguese for (film) director, then the Portuguese for to film.
\end{quote}

\subsection*{You'll want to know: o guião}
\textbf{o guião} — \textquotedblleft{}script, screenplay\textquotedblright{}.

\begin{quote}
O guião do filme é baseado num romance.
\end{quote}

\begin{quote}
\emph{The film's script is based on a novel.}
\end{quote}

In Brazil: \emph{o roteiro}.

\begin{cousinweb}
\emph{guião} is linked in EtymDB with \emph{guia} \textquotedblleft{}a guide\textquotedblright{}: the text that guides the film. The verb \emph{guiar} is usually traced to a Germanic word.
\end{cousinweb}

\begin{grammarlens}[title={What you've built}]
One more word for telling, explaining and describing at B1.
\end{grammarlens}

\subsection*{Your turn}
\begin{itemize}
\raggedright
  \item \emph{Say it:} \emph{o guião}
  \item \emph{Say it:} \emph{o guião}, once more
  \item \emph{Say it:} a sentence of your own with \emph{o guião}
  \item \emph{Recall:} say the Portuguese for (film) director, then the Portuguese for to film, then say \emph{o guião} again
\end{itemize}

\begin{tcolorbox}[breakable,colback=teal!4,colframe=teal!35!black,title={Before you move on}]
What does o guião mean? (\textquotedblleft{}Script, screenplay\textquotedblright{}.) Say it once more.
\end{tcolorbox}
\section[o documentário]{o documentário — documentary}

\label{lesson:PT-C527-documentario}



\begin{quote}
Before the new one: say the Portuguese for to film, then the Portuguese for script.
\end{quote}

\subsection*{You'll want to know: o documentário}
\textbf{o documentário} — \textquotedblleft{}documentary\textquotedblright{}.

\begin{quote}
Vimos um documentário sobre os oceanos.
\end{quote}

\begin{quote}
\emph{We watched a documentary about the oceans.}
\end{quote}

\emph{Um documentário sobre} — a documentary about.

\begin{cousinweb}
\emph{documentário} is built on \emph{documento} \textquotedblleft{}a document\textquotedblright{}, from Latin \emph{documentum} \textquotedblleft{}a lesson, a proof\textquotedblright{} (Lewis \& Short), from \emph{docēre} \textquotedblleft{}to teach\textquotedblright{}.
\end{cousinweb}

\begin{grammarlens}[title={What you've built}]
One more word for telling, explaining and describing at B1.
\end{grammarlens}

\subsection*{Your turn}
\begin{itemize}
\raggedright
  \item \emph{Say it:} \emph{o documentário}
  \item \emph{Say it:} \emph{o documentário}, once more
  \item \emph{Say it:} a sentence of your own with \emph{o documentário}
  \item \emph{Recall:} say the Portuguese for to film, then the Portuguese for script, then say \emph{o documentário} again
\end{itemize}

\begin{tcolorbox}[breakable,colback=teal!4,colframe=teal!35!black,title={Before you move on}]
What does o documentário mean? (\textquotedblleft{}Documentary\textquotedblright{}.) Say it once more.
\end{tcolorbox}
\section[o elenco]{o elenco — cast (of a film or play)}

\label{lesson:PT-C527-elenco}



\begin{quote}
Before the new one: say the Portuguese for script, then the Portuguese for documentary.
\end{quote}

\subsection*{You'll want to know: o elenco}
\textbf{o elenco} — \textquotedblleft{}cast (of a film or play)\textquotedblright{}.

\begin{quote}
O elenco do filme é quase todo português.
\end{quote}

\begin{quote}
\emph{The cast of the film is almost all Portuguese.}
\end{quote}

Also a list or a squad: \emph{o elenco da equipa}.

\begin{cousinweb}
\emph{elenco} comes from Late Latin \emph{elenchus} (EtymDB), from Greek \emph{élenkhos} \textquotedblleft{}a scrutiny, a test\textquotedblright{}; the sense \textquotedblleft{}a list\textquotedblright{} grew up later.
\end{cousinweb}

\begin{grammarlens}[title={What you've built}]
One more word for telling, explaining and describing at B1.
\end{grammarlens}

\subsection*{Your turn}
\begin{itemize}
\raggedright
  \item \emph{Say it:} \emph{o elenco}
  \item \emph{Say it:} \emph{o elenco}, once more
  \item \emph{Say it:} a sentence of your own with \emph{o elenco}
  \item \emph{Recall:} say the Portuguese for script, then the Portuguese for documentary, then say \emph{o elenco} again
\end{itemize}

\begin{tcolorbox}[breakable,colback=teal!4,colframe=teal!35!black,title={Before you move on}]
What does o elenco mean? (\textquotedblleft{}Cast (of a film or play)\textquotedblright{}.) Say it once more.
\end{tcolorbox}
